\chapter{Finite \texorpdfstring{\(q\)}{q}-Pochhammer inversion}
\label{chap:finite-q-pochhammer-inversion}

\section{The finite product as an inverse covering}
\label{fp-finite-covering}

Fix \(n\geq1\) and write
\[
 P_n(a,q):=\QPochhammer{a}{q}{n}=\prod_{j=0}^{n-1}(1-aq^j).
\]
Throughout, \(\QInteger{n}{q}\) and \(\GaussianBinomial{n}{r}{q}\) denote their polynomial
continuations at \(q=0,1\) and at roots of unity; no quotient with a vanishing
denominator is being used there.
This section consolidates the finite-product results that were repeated in the
six source reports.  The results are human-proved frontier results: no Lean
formalization is claimed here.  We first invert in \(a\), then in \(q\), and
keep the two branch geometries separate.

Whenever no factor vanishes, logarithmic differentiation gives
\begin{align}
 \partial_a\log P_n(a,q)
  &=-\sum_{j=0}^{n-1}\frac{q^j}{1-aq^j},
  \label{fp-a-logder}\\
 \partial_q\log P_n(a,q)
  &=-a\sum_{j=1}^{n-1}\frac{j q^{j-1}}{1-aq^j}.
  \label{fp-q-logder}
\end{align}
Both identities also hold after clearing denominators at a zero.  They will be
used only on zero-free sets, where the logarithms are unambiguous locally.

\subsection{Argument inversion at the unit level}

Euler's finite \(q\)-binomial identity is
\begin{equation}
 P_n(a,q)=\sum_{r=0}^{n}(-1)^r q^{\binom r2}
 \GaussianBinomial{n}{r}{q}a^r.
 \label{fp-finite-q-binomial}
\end{equation}
Set \(s=1-y\) and
\[
 b_r=(-1)^{r+1}q^{\binom r2}\GaussianBinomial{n}{r}{q},
 \qquad f(a)=1-P_n(a,q)=\sum_{r\geq1}b_ra^r.
\]

\begin{theorem}[All-order argument inverse at \(y=1\)]
\label{fp-unit-inverse}
If \(\QInteger{n}{q}=b_1\ne0\), there is a unique analytic branch \(a=a(s)\) with
\(a(0)=0\) and \(P_n(a(s),q)=1-s\).  Its coefficients are
\begin{equation}
 [s^m]a(s)=\frac1m[z^{m-1}]
 \left(\frac{z}{1-P_n(z,q)}\right)^m,
 \qquad m\geq1.
 \label{fp-unit-lagrange}
\end{equation}
In particular,
\begin{align}
 a={}&\frac{s}{\QInteger{n}{q}}
 +\frac{q\GaussianBinomial{n}{2}{q}}{\QInteger{n}{q}^3}s^2 \notag\\
 &+\left(
 \frac{2q^2\GaussianBinomial{n}{2}{q}^{\,2}}{\QInteger{n}{q}^5}
 -\frac{q^3\GaussianBinomial{n}{3}{q}}{\QInteger{n}{q}^4}
 \right)s^3+\BigO(s^4).
 \label{fp-unit-first-terms}
\end{align}
\end{theorem}

\begin{proof}
The linear coefficient of \(f\) is \(f'(0)=\QInteger{n}{q}\ne0\), so the analytic
inverse-function theorem supplies one and only one local branch.  Rewrite
\(s=f(a)\) as \(a=s\phi(a)\), where \(\phi(a)=a/f(a)\).  Lagrange inversion
gives \([s^m]a=m^{-1}[z^{m-1}]\phi(z)^m\), which is
\cref{fp-unit-lagrange}.  Finally
\[
 s=\QInteger{n}{q}a-q\GaussianBinomial{n}{2}{q}a^2
   +q^3\GaussianBinomial{n}{3}{q}a^3+\BigO(a^4).
\]
Substitute \(a=A_1s+A_2s^2+A_3s^3+\cdots\) and compare coefficients.  This
gives
\[
 A_1=\QInteger{n}{q}^{-1},\quad
 A_2=\frac{q\GaussianBinomial{n}{2}{q}}{\QInteger{n}{q}^3},\quad
 A_3=\frac{2q^2\GaussianBinomial{n}{2}{q}^{\,2}}{\QInteger{n}{q}^5}
       -\frac{q^3\GaussianBinomial{n}{3}{q}}{\QInteger{n}{q}^4},
\]
proving both displayed claims.
\end{proof}

\subsection{All zero sheets and the complete real atlas}

For \(q\ne0\), the argument zeros are \(a_j=q^{-j}\), \(0\leq j<n\).
They are distinct precisely when \(q^r\ne1\) for \(1\leq r<n\).

\begin{proposition}[Exact normalization at every simple zero]
\label{fp-zero-normalization}
Assume \(q\ne0\) and \(q^r\ne1\) for \(1\leq r<n\).  Then
\begin{equation}
 \partial_aP_n(q^{-j},q)=
 (-1)^{j+1}q^{-j(j-1)/2}
 \QPochhammer{q}{q}{j}\QPochhammer{q}{q}{n-1-j}.
 \label{fp-zero-derivative}
\end{equation}
Consequently the inverse sheet through \((y,a)=(0,q^{-j})\) begins
\begin{equation}
 a(y)=q^{-j}+
 \frac{(-1)^{j+1}q^{j(j-1)/2}}
 {\QPochhammer{q}{q}{j}\QPochhammer{q}{q}{n-1-j}}\,y+\BigO(y^2).
 \label{fp-zero-inverse}
\end{equation}
\end{proposition}

\begin{proof}
Differentiate the unique vanishing factor:
\[
 \partial_aP_n(q^{-j},q)
 =-q^j\prod_{\ell\ne j}(1-q^{\ell-j}).
\]
For \(\ell<j\), write
\(1-q^{\ell-j}=-q^{-(j-\ell)}(1-q^{j-\ell})\).  Their product is
\((-1)^jq^{-j(j+1)/2}\QPochhammer{q}{q}{j}\), while the factors with
\(\ell>j\) give \(\QPochhammer{q}{q}{n-1-j}\).  Multiplication by \(-q^j\) proves
\cref{fp-zero-derivative}.  Its nonvanishing permits ordinary analytic
inversion, and division by the derivative gives \cref{fp-zero-inverse}.
\end{proof}

\begin{theorem}[Complete real argument atlas]
\label{fp-real-atlas}
Let \(q>0\), \(q\ne1\), and order the zeros increasingly as
\(r_0<\cdots<r_{n-1}\).  Then:
\begin{enumerate}
\item \(P_n(\,\cdot\,,q)\) is positive and strictly decreasing on
  \((-\infty,r_0)\), and maps that interval bijectively onto
  \((0,\infty)\).
\item Every gap \((r_j,r_{j+1})\) contains exactly one critical point
  \(c_j\).  It is simple, so the two inverse sheets meeting at
  \(P_n(c_j,q)\) have square-root ramification.
\item There is no critical point on \((r_{n-1},\infty)\).
\end{enumerate}
For \(n\ge2\), the maximal real intervals on which
\(P_n(\,\cdot\,,q)\) is one-to-one are
\[
 (-\infty,c_0],\qquad
 [c_{j-1},c_j]\quad(1\leq j\leq n-2),\qquad
 [c_{n-2},\infty).
\]
Each contains exactly one simple zero \(r_j\), and together they exhaust the
maximal real injective branches; their interiors are the maximal open domains
of regular real inversion.  Equivalently, each zero-free gap
\((r_j,r_{j+1})\) consists of two half-branches meeting at the fold \(c_j\).
When \(n=1\), the sole maximal branch is all of \(\RealNumbers\).
\end{theorem}

\begin{proof}
Since \(1-aq^j=-q^j(a-q^{-j})\), away from the roots
\begin{equation}
 \frac{\partial_aP_n(a,q)}{P_n(a,q)}
 =\sum_{j=0}^{n-1}\frac1{a-r_j},
 \qquad
 \frac{\Differential}{\Differential a}\frac{\partial_aP_n}{P_n}
 =-\sum_{j=0}^{n-1}\frac1{(a-r_j)^2}<0.
 \label{fp-root-logder}
\end{equation}
On the leftmost interval every summand in the first expression is negative.
All product factors are positive there, the product tends to \(+\infty\) at
\(-\infty\), and it tends to zero at \(r_0\); this proves the first item.
On \((r_j,r_{j+1})\), the logarithmic derivative decreases continuously from
\(+\infty\) to \(-\infty\), so it has exactly one zero.  At that zero,
\(P_n''/P_n\) equals the strictly negative second expression in
\cref{fp-root-logder}; hence \(P_n''\ne0\).  Taylor's formula and reversion of
its first nonzero quadratic term give two square-root sheets.  To the right of
the largest root every summand in \cref{fp-root-logder} is positive.  The
preceding argument, together with the nonvanishing derivatives at the simple
zeros, proves that \(c_0,\ldots,c_{n-2}\) are all the real zeros of
\(\partial_aP_n\).  Its sign is therefore fixed and nonzero on each component
of their complement.  The function is strictly monotone on the closure of
each such component.  Because every \(c_j\) is a strict local extremum, no
closure can be extended past one of its finite endpoints while preserving
injectivity.  These closures are therefore precisely the maximal injective
intervals, while their interiors are the maximal regular inverse domains.
The interlacing
\(c_{j-1}<r_j<c_j\), with the evident one-sided versions at the extremes,
shows that each contains exactly one simple zero; analytic inversion continues
through that zero.  For \(n=1\), the polynomial is linear and the same
conclusion is immediate.
\end{proof}

Put \(N=\binom n2\), \(A_n=(-1)^nq^N\), and
\(e_1=\sum_{j=0}^{n-1}q^{-j}\).  Factoring the leading power of \(a\) gives
\[
 P_n(a,q)=A_na^n\left(1-\frac{e_1}{a}+\BigO(a^{-2})\right).
\]
Thus, on each choice \(T=(y/A_n)^{1/n}\),
\begin{equation}
 a=T+\frac{e_1}{n}+\BigO(T^{-1}).
 \label{fp-argument-infinity}
\end{equation}
Indeed, substitution of \(a=T+c+\BigO(T^{-1})\) makes the coefficient of
\(T^{-1}\) equal to \(nc-e_1\), which forces \(c=e_1/n\).

\subsection{Cyclotomic collisions in the argument}

\begin{theorem}[Exact root-of-unity factorization and ramification]
\label{fp-cyclotomic-argument}
Let \(\zeta\) be primitive of order \(d\), and write \(n=Md+r\) with
\(0\leq r<d\).  Then
\begin{equation}
 \QPochhammer{a}{\zeta}{n}=(1-a^d)^M\QPochhammer{a}{\zeta}{r}.
 \label{fp-root-block}
\end{equation}
At \(a_s=\zeta^{-s}\), \(0\leq s<d\), the zero multiplicity is
\begin{equation}
 \nu_s=M+\IndicatorOf{s<r}.
 \label{fp-root-multiplicity}
\end{equation}
For every \(s\) with \(\nu_s\ge1\), the point \(a_s\) is a zero of
multiplicity \(\nu_s\).  On the ramified cover of a sufficiently small
punctured target neighborhood, the inverse at target zero then has exactly
\(\nu_s\) local Puiseux sheets coalescing at \(a_s\).  Equivalently, after
choosing a \(\nu_s\)-th root on a simply connected punctured target domain,
these are the \(\nu_s\) determinations obtained by multiplying that root by
the \(\nu_s\)-th roots of unity.  When \(\nu_s=0\), \(a_s\) is not a zero
and contributes no inverse germ at target zero.
At the unit target \(y=1\), the branch through \(a=0\) is regular when
\(r>0\).  When \(r=0\), it has the \(d\) local sheets
\begin{equation}
 a=\omega\left(\frac{1-y}{M}\right)^{1/d}
   +\BigO((1-y)^{2/d}),\qquad \omega^d=1.
 \label{fp-unit-cyclotomic-puiseux}
\end{equation}
\end{theorem}

\begin{proof}
Partition the indices into \(M\) complete residue blocks modulo \(d\) and one
residual block.  Each complete block contributes
\(\prod_{j=0}^{d-1}(1-a\zeta^j)=1-a^d\), proving
\cref{fp-root-block}.  A factor vanishes at \(a_s\) exactly when its index is
congruent to \(s\pmod d\).  That residue occurs \(M\) times in the complete
blocks and once more precisely when \(s<r\), proving
\cref{fp-root-multiplicity}.  When \(\nu_s\ge1\), removing the nonvanishing
factors leaves a nonzero constant times \((a-a_s)^{\nu_s}\), whose equation
with \(y\) has \(\nu_s\) Puiseux determinations.  When \(\nu_s=0\), no factor
vanishes at \(a_s\), so there is no inverse germ over target zero there.

At \(a=0\), the derivative is
\(-\sum_{j=0}^{n-1}\zeta^j=-\sum_{j=0}^{r-1}\zeta^j\), which is nonzero for
\(r>0\).  If \(r=0\), \cref{fp-root-block} gives
\(\QPochhammer{a}{\zeta}{n}=(1-a^d)^M=1-Ma^d+\BigO(a^{2d})\); solving its leading equation
gives \cref{fp-unit-cyclotomic-puiseux}.
\end{proof}

At \(q=-1\), this includes
\[
 \QPochhammer{a}{-1}{2M}=(1-a^2)^M,\qquad
 \QPochhammer{a}{-1}{2M+1}=(1-a)(1-a^2)^M.
\]
For even order every argument sheet is obtained from
\(a=\pm\sqrt{1-\eta y^{1/M}}\), where \(\eta^M=1\).

\section{Finite-product inversion in the base}
\label{fp-base-inversion}

\subsection{The canonical real base branch}

\begin{theorem}[Global finite base inverse on \(0\leq q\leq1\)]
\label{fp-global-base}
Let \(n\geq2\).
\begin{enumerate}
\item If \(0<a<1\), then \(q\mapsto P_n(a,q)\) is strictly decreasing from
  \(1-a\) to \((1-a)^n\).  Every
  \((1-a)^n<y<1-a\) has exactly one inverse \(q\in(0,1)\).
\item If \(a<0\), the map is strictly increasing from \(1-a\) to
  \((1-a)^n\).  Every \(1-a<y<(1-a)^n\) has exactly one inverse
  \(q\in(0,1)\).
\end{enumerate}
For \(a=0\), \(a=1\), or \(n=1\), the map is constant in \(q\).
\end{theorem}

\begin{proof}
For \(0<q<1\), every denominator in \cref{fp-q-logder} is positive.  Its sign
is therefore \(-\Sign(a)\), and the nonzero derivative proves strict
monotonicity.  Direct substitution gives \(P_n(a,0)=1-a\) and
\(P_n(a,1)=(1-a)^n\).  Continuity and strict monotonicity now give the stated
ranges and uniqueness.  The three constant cases follow immediately.
\end{proof}

\subsection{\texorpdfstring{All-order inversion at \(q=0\)}{All-order inversion at q=0}}

Assume \(n\geq2\) and \(a\notin\{0,1\}\), and normalize the target by
\begin{equation}
 \eta=\frac{(1-a)-y}{a(1-a)}.
 \label{fp-small-base-coordinate}
\end{equation}
Define \(F_{n,a}(q)=(1-P_n(a,q)/(1-a))/a\).  Then
\(F_{n,a}(q)=q+\BigO(q^2)\), and \(P_n(a,q)=y\) is precisely
\(F_{n,a}(q)=\eta\).

\begin{theorem}[Exact coefficient engine at \(q=0\)]
\label{fp-small-base-engine}
The inverse branch \(q=Q_{n,a}(\eta)\) is analytic at zero and satisfies
\begin{equation}
 [\eta^m]Q_{n,a}(\eta)=\frac1m[z^{m-1}]
 \left(\frac{z}{F_{n,a}(z)}\right)^m.
 \label{fp-small-base-lagrange}
\end{equation}
Moreover,
\begin{equation}
 \log\frac{P_n(a,q)}{1-a}=\sum_{r\geq1}\ell_r^{(n)}(a)q^r,\qquad
 \ell_r^{(n)}(a)=
 -\sum_{\substack{d\mid r\\1\leq d\leq n-1}}
 \frac{a^{r/d}}{r/d}.
 \label{fp-small-base-divisor}
\end{equation}
For \(n\geq6\),
\begin{equation}
 q=\eta-\eta^2+(1+a)\eta^3-(1+4a)\eta^4
 +(1+11a+3a^2)\eta^5+\BigO(\eta^6).
 \label{fp-small-base-five}
\end{equation}
For \(n=2\), \(q=\eta\) exactly; for \(n=3\), the cubic coefficient is
\(2+a\); and for every \(n\geq4\) it is \(1+a\).
\end{theorem}

\begin{proof}
The coefficient of \(q\) in \(F_{n,a}\) is one, so analytic inversion and
Lagrange's formula give \cref{fp-small-base-lagrange}.  For the divisor
formula, expand
\[
 \log\frac{P_n(a,q)}{1-a}
 =-\sum_{j=1}^{n-1}\sum_{m\geq1}\frac{a^mq^{jm}}m.
\]
The coefficient of \(q^r\) comes from \(jm=r\), proving
\cref{fp-small-base-divisor}.

When \(n\geq6\), subset expansion gives
\[
 F_{n,a}(q)=q+q^2+(1-a)q^3+(1-a)q^4+(1-2a)q^5+\BigO(q^6).
\]
Substitute \(q=\eta+A_2\eta^2+\cdots+A_5\eta^5\) and compare powers.
Successively,
\(A_2=-1\), \(A_3=1+a\), \(A_4=-(1+4a)\), and
\(A_5=1+11a+3a^2\).  The same calculation with the available positive
indices gives the stated low-order exceptions.
\end{proof}

\subsection{\texorpdfstring{The complete logarithmic jet at \(q=1\)}{The complete logarithmic jet at q=1}}

Put \(q=\EulerE^{-t}\), \(S_m(n)=\sum_{j=0}^{n-1}j^m\), and choose the local
logarithm of \(y/(1-a)^n\).

\begin{theorem}[All-order regular base inverse at \(q=1\)]
\label{fp-unit-base-inverse}
Let \(n\geq2\) and \(a\in\ComplexNumbers\setminus\{0,1\}\).  With
\[
 u=\log\frac{y}{(1-a)^n},\qquad
 c_m=\frac{(-1)^{m+1}}{m!}\Polylogarithm_{1-m}(a)S_m(n),
\]
the forward germ and its inverse are
\begin{align}
 u&=\sum_{m\geq1}c_mt^m,
 \label{fp-unit-base-forward}\\
 t&=\frac{u}{c_1}-\frac{c_2}{c_1^3}u^2
 +\left(\frac{2c_2^2}{c_1^5}-\frac{c_3}{c_1^4}\right)u^3
 +\BigO(u^4),\qquad q=\EulerE^{-t}.
 \label{fp-unit-base-reverse}
\end{align}
Here \(c_1=a\binom n2/(1-a)\ne0\), and all subsequent inverse coefficients
are obtained uniquely by triangular coefficient comparison.
\end{theorem}

\begin{proof}
For \(|a|<1\),
\[
 \log(1-a\EulerE^{-x})
 =-\sum_{r\geq1}\frac{a^r\EulerE^{-rx}}r
 =\log(1-a)+\sum_{m\geq1}
   \frac{(-1)^{m+1}}{m!}\Polylogarithm_{1-m}(a)x^m.
\]
Substitute \(x=jt\) and sum over \(0\leq j<n\).  Both sides are analytic in
\(a\) away from \(a=1\), and the nonpositive-index polylogarithms are rational,
so analytic continuation proves the formula for every stated \(a\).
Since \(S_1(n)=\binom n2\), the linear coefficient is nonzero.  Substitution
of \(t=A_1u+A_2u^2+A_3u^3+\cdots\) gives
\(A_1=c_1^{-1}\), \(A_2=-c_2c_1^{-3}\), and
\(A_3=2c_2^2c_1^{-5}-c_3c_1^{-4}\).  Every later equation is linear in the
new coefficient.
\end{proof}

For \(n\geq2\), holding \(y\ne0\) fixed while \(q\to1\) gives a useful mixed
inverse.  Choose \(r=y^{1/n}\), put \(a_0=1-r\), and keep that root
determination fixed.
Substitution in the logarithm of the level equation gives
\begin{equation}
 a(t)=a_0+\frac{n-1}{2}a_0t+
 \frac{a_0(n-1)\bigl(2n-4-3a_0(n-1)\bigr)}
 {24(1-a_0)}t^2+\BigO(t^3).
 \label{fp-fixed-target-argument}
\end{equation}
Indeed, writing \(a=a_0+At+Bt^2\) makes the constant term
\(n\log(1-a_0)=\log y\); the linear and quadratic equations give
\(A=(n-1)a_0/2\) and the displayed \(B\).

\subsection{\texorpdfstring{The exceptional slice \(q=-1\)}{The exceptional slice q=-1}}
\label{fp-minus-one-section}

\begin{theorem}[Complete zero-free criticality at \(q=-1\)]
\label{fp-minus-one-criticality}
Let \(n\geq3\) and \(a\notin\{0,\pm1\}\).  Then
\(\partial_qP_n(a,-1)=0\) if and only if
\begin{equation}
 a=a_*(n)=
 \begin{cases}
  1/(n-1),&n\text{ even},\\
  -1/n,&n\text{ odd}.
 \end{cases}
 \label{fp-minus-one-parameter}
\end{equation}
For \(n=2s\), \(s\geq2\),
\begin{align}
 y_*&=\left(\frac{4s(s-1)}{(2s-1)^2}\right)^s,
 \label{fp-minus-one-even-value}\\
 L_2&:=\partial_q^2\log P_{2s}(a_*,-1)
 =-\frac{(s+1)(2s-1)^2}{12s(s-1)}.
 \label{fp-minus-one-even-second}
\end{align}
For \(n=2s+1\), \(s\geq2\),
\begin{align}
 y_*&=\left(\frac{2s+2}{2s+1}\right)^{s+1}
      \left(\frac{2s}{2s+1}\right)^s,
 \label{fp-minus-one-odd-value}\\
 L_2&=-\frac{(s-1)(2s+1)^2}{12s(s+1)}.
 \label{fp-minus-one-odd-second}
\end{align}
All these slices are quadratic folds, with local inverse
\begin{equation}
 q=-1\pm\sqrt{\frac{2(y-y_*)}{y_*L_2}}+\BigO(y-y_*).
 \label{fp-minus-one-fold}
\end{equation}
Within this zero-free family, the sole higher degeneracy is
\begin{equation}
 P_3(-1/3,q)=\frac{32}{27}+\frac4{27}(q+1)^3.
 \label{fp-minus-one-cubic}
\end{equation}
\end{theorem}

\begin{proof}
Let \(L=\log P_n\).  At \(q=-1\), separation of even and odd indices in
\cref{fp-q-logder} gives
\[
 L_q(a,-1)=a\left(\frac{E}{1-a}-\frac{O}{1+a}\right),
\]
where \(E\) and \(O\) are the sums of the positive even and odd indices below
\(n\).  For \(n=2s\), \(E=s(s-1)\), \(O=s^2\), giving
\(a=1/(2s-1)\).  For \(n=2s+1\), \(E=s(s+1)\), \(O=s^2\), giving
\(a=-1/(2s+1)\).  Substitution in
\(P_n(a,-1)=(1-a)^{\Ceiling{n/2}}(1+a)^{\Floor{n/2}}\) gives the
two critical values.

Differentiating \cref{fp-q-logder} once more gives
\[
 L_{qq}=-a\sum_{j=1}^{n-1}\left(
 \frac{j(j-1)q^{j-2}}{1-aq^j}
 +\frac{aj^2q^{2j-2}}{(1-aq^j)^2}\right).
\]
For completeness, the four parity sums needed here are
\[
\begin{array}{c|cc}
 &\displaystyle\sum j(j-1)&\displaystyle\sum j^2\\ \hline
 2,4,\ldots,2(s-1)&
 \dfrac{s(s-1)(4s-5)}3&
 \dfrac{2s(s-1)(2s-1)}3\\[0.7em]
 1,3,\ldots,2s-1&
 \dfrac{s(s-1)(4s+1)}3&
 \dfrac{s(4s^2-1)}3\\[0.7em]
 2,4,\ldots,2s&
 \dfrac{s(s+1)(4s-1)}3&
 \dfrac{2s(s+1)(2s+1)}3 .
\end{array}
\]
At \(q=-1\), the first summand in \(L_{qq}\) uses the even row with a plus
sign and the odd row with a minus sign, whereas its \(j^2\) summand uses both
rows with a plus sign.  For \(n=2s\), insert the first two rows and
\(a=1/(2s-1)\); for \(n=2s+1\), insert the last two rows and
\(a=-1/(2s+1)\).  Cancelling the common factors gives exactly
\cref{fp-minus-one-even-second,fp-minus-one-odd-second}.  The even expression
never vanishes for \(s\geq2\); the odd expression vanishes only at \(s=1\),
where direct multiplication proves \cref{fp-minus-one-cubic}.
At a critical point \(P_{qq}=P_nL_{qq}=y_*L_2\ne0\).  Taylor expansion and
reversion of the quadratic term give \cref{fp-minus-one-fold}.
\end{proof}

The exclusions in \cref{fp-minus-one-criticality} are structural, not hidden
exceptions.  At \(a=0\) the product is identically one, and at \(a=1\) it is
identically zero.  At \(a=-1\),
\[
 P_n(-1,q)=\prod_{j=0}^{n-1}(1+q^j)
\]
has at \(q=-1\) a zero of exact order \(r=\Floor{n/2}\), because
exactly the odd-indexed factors vanish and each does so simply.  More
precisely, with \(e=\Ceiling{n/2}\) and
\(J_n=\prod_{\substack{1\leq j<n\\j\ {\rm odd}}}j\),
\begin{equation}
 P_n(-1,q)=2^eJ_n(q+1)^r+\BigO((q+1)^{r+1}).
 \label{fp-minus-one-zero-slice}
\end{equation}
Indeed, every even-indexed factor tends to two and
\(1+q^j=j(q+1)+\BigO((q+1)^2)\) for odd \(j\).  Thus, on the ramified target
cover,
\[
 q=-1+\omega\left(\frac{y}{2^eJ_n}\right)^{1/r}
 +\BigO(y^{2/r}),\qquad \omega^r=1.
\]
This zero-target cyclotomic ramification is the repeated-factor phenomenon
seen in \cref{fp-cyclotomic-argument}; it is different from the nonzero-level
critical points classified above.

\begin{corollary}[The cubic slice lies on a Morse saddle]
\label{fp-minus-one-morse}
For every odd \(n=2s+1\), \(s\geq1\), the point
\((a,q)=(-1/(2s+1),-1)\) is a nondegenerate saddle of
\((a,q)\mapsto P_n(a,q)\).  Thus \cref{fp-minus-one-cubic} is a tangent
one-dimensional slice, not a cusp of the full parameter surface.
\end{corollary}

\begin{proof}
At this point both \(L_a\) and \(L_q\) vanish.  Direct differentiation of
\(L=\sum_{j=0}^{2s}\log(1-aq^j)\), followed by the parity sums, gives
\begin{align*}
 L_{aa}&=-\frac{(2s+1)^3}{4s(s+1)},&
 L_{aq}&=-\frac{(2s+1)^2}{4(s+1)},&
 L_{qq}&=-\frac{(s-1)(2s+1)^2}{12s(s+1)}.
\end{align*}
For instance,
\[
 L_{aa}=-\frac{s+1}{(1-a)^2}-\frac{s}{(1+a)^2},\qquad
 L_{aq}=\frac{s(s+1)}{(1-a)^2}-\frac{s^2}{(1+a)^2},
\]
which yield the first two entries after substituting \(a=-1/(2s+1)\); the
third is \cref{fp-minus-one-odd-second}.  Hence
\begin{equation}
 \det\operatorname{Hess}L
 =-\frac{(2s+1)^4(s^2+s+1)}{48s^2(s+1)^2}<0.
 \label{fp-minus-one-hessian}
\end{equation}
Since \(P_n\ne0\) and its gradient vanishes, its Hessian is \(P_n\) times that
of \(L\).  It is therefore also nonsingular and indefinite.
\end{proof}

\subsection{Reciprocal-base inversion at infinity}

Let \(n\geq2\), \(a\notin\{0,1\}\), \(N=\binom n2\),
\(C=(1-a)(-a)^{n-1}\), and
\[
 H(z)=\prod_{j=1}^{n-1}(1-a^{-1}z^j).
\]
The exact factorization \(P_n(a,q)=Cq^NH(q^{-1})\) converts infinity into an
ordinary point.

\begin{theorem}[All-order finite base inverse at infinity]
\label{fp-base-infinity}
Choose an \(N\)-th root \(x=(C/y)^{1/N}\).  The corresponding branch
\(s=q^{-1}\) is the convergent solution of
\begin{equation}
 s=xH(s)^{1/N},\qquad
 [x^m]s(x)=\frac1m[z^{m-1}]H(z)^{m/N}.
 \label{fp-infinity-lagrange}
\end{equation}
In particular,
\begin{equation}
 q=x^{-1}+\frac1{aN}+\BigO(x).
 \label{fp-infinity-first}
\end{equation}
All \(N\) complex sheets arise from the \(N\) choices of \(x\).
\end{theorem}

\begin{proof}
The level equation is \(y=Cs^{-N}H(s)\), hence
\(x=sH(s)^{-1/N}\), equivalently the first equation in
\cref{fp-infinity-lagrange}.  Since \(H(0)=1\), a consistent analytic
\(N\)-th root exists near zero and the analytic inverse-function theorem gives
a unique convergent branch.  Lagrange inversion proves the coefficient
formula.  Finally \(H(s)=1-a^{-1}s+\BigO(s^2)\), so
\(s=x-a^{-1}x^2/N+\BigO(x^3)\).  Taking the reciprocal gives
\cref{fp-infinity-first}.
\end{proof}

\begin{warning}[Claims deliberately not promoted]
The source reports contain exact-looking secondary-discriminant factorizations
for a proposed first Maxwell collision at \(n=5\).  No independent exact
certificate for the large factorization is currently part of the canonical
package, so that special assertion is not stated here as a theorem.  The
proved branch atlas remains valid, and a certified resultant/Sturm computation
would suffice to add the collision theorem later.
\end{warning}
